% C05 continuation for ED26-BM-RAT2021-15-1. New reader material only.
% Root-owned placement: after the core's explicit-rationality/argumentation
% discussion. Existing core, prior evidence fragment and P11 record unchanged.
\subsection{Collaborative selection is not yet developmental transfer}
\label{sec:p19-r151-collaborative-interface}

Moshman and Geil's 1998 original experiment already has its full methodological
and numerical record in \cref{sec:p11-ri-individual-social}, with response-unit
conventions at \cref{def:p11-ri-collaborative-statistics}. It is the same
experiment cited in the developmental chapter, not a second independent
demonstration. Course-based undergraduates completed an abstract four-card
task individually or in consensus-oriented groups of five or six, with
written explanations. The original source gives 143 participants but compares
32 individual outputs with 20 group outputs. Its pooled 15/20 correct group
solutions and 3/32 correct individual solutions are not an age comparison
or 143 independent group trials
\cite[pp.~233--236]{MoshmanGeil1998CollaborativeReasoning}.

The developmental relevance lies in a possible mechanism: participants can
combine partial insights into a justified argument. In the two illustrated
groups, an argument for checking the odd-number card was combined with
an argument against checking the even-number card; neither group initially
contained a member choosing the complete correct pattern. Such construction
is different from counting votes or copying an initially complete answer.
Those illustrations were purposively selected from three all-initially-
incorrect groups; the third videotape was lost. The paper supplies no
formal process-coding reliability or isolated mediation estimate
\cite[pp.~237--245]{MoshmanGeil1998CollaborativeReasoning}.

The original discussion nevertheless does not establish a general childhood
developmental effect. Participant ages and a child cohort are not supplied.
Responses are immediate, on the same task, following discussion; no delayed
or novel-task test demonstrates lasting individual transfer. Group discussions
took approximately two-and-a-half times the individual solving time, without
a time-yoked control. Eliminating simple adoption of the initial majority
answer does not eliminate every influence of confidence, status or later
consensus. These limits accompany the existing original-study record rather
than being replaced by a generic claim that collaboration improves rationality.

At p.~246 the authors explicitly decline to claim durable, generalizable
competence change from this brief interaction. Their prediction that groups
of children younger than eleven or twelve might fare no better than individual
children, if the task requires a particular metalogical competence, is a
conditional prediction, not an observed child result. Repeated dialogic
learning is a proposed developmental route. The two adolescent interventions
below test related but different outcomes; they cannot retroactively turn
the undergraduate selection experiment into childhood longitudinal evidence
\cite[pp.~242--246]{MoshmanGeil1998CollaborativeReasoning}.

\subsection{Kuhn and Udell's 2003 intervention: discourse and its individual product}
\label{sec:p19-r151-kuhn-udell}

Kuhn and Udell distinguish an individual argument supporting a claim from
argumentative discourse responding to an interlocutor. Their intervention
targets a shift from articulating one's own position to engaging the
opponent's reasons, including critique and rebuttal. The assessment criterion
is a source-defined discourse/argument rubric, not an answer key identifying
the ethically correct position on capital punishment. More coded
counterarguments do not by themselves prove greater logical validity
\cite[pp.~1245--1248]{KuhnUdell2003Development}.

The analyzed sample comprised 34 eighth graders from two low-performing
New York City public middle schools. The article describes its target group
as academically at-risk thirteen-to-fourteen-year-olds but gives no individual
age distribution. It reports 14 Hispanic, 19 African-American and one
Ethiopian participant; 11 male and 23 female. These are the reported
sampling descriptors, not diagnoses or population-wide competence judgments.
There were 21 completers in the experimental condition and 13 in comparison.
Five additional starters, three experimental and two comparison, did not
complete because of chronic absence or suspension. Thus 39 started and
34 completed; outcome-specific samples are smaller in some analyses.

The authors report random assignment to conditions without specifying its
algorithm. Teams were then formed from pro/con opinions. The experimental
teams were represented at both schools; the comparison teams were at one
school. Team sizes were four to eight, with typical attendance three to
six. Recruitment/response rates, consent details and native-language
eligibility are not reported. English task formulations are supplied.
School, team and discussion-partner dependence remain relevant to inference
\cite[pp.~1247--1250]{KuhnUdell2003Development}.

The full condition received sixteen sessions scheduled twice weekly for
ninety minutes. Interruptions extended the nominal eight-week schedule to
approximately twelve weeks. Both conditions shared the first seven sessions,
developing reasons, elaborating them, using evidence, evaluating reasons,
and assembling an argument. Comparison activity ended after this phase;
it did not receive nine matched-duration alternative sessions. The full
condition continued with opposing-side reasons, counterarguments, rebuttals,
mixed evidence, practice dialogues and a final debate. The difference is
therefore a scaffolded package with more exposure, not a dialogue-only
manipulation at equal dose.
Comparison pupils also engaged in same-side team activity; this was not
an individual-only or no-collaboration control.

Adult coaches supported the activities. Participants made reason cards,
considered how an opponent would react, attached evidence and rebuttal cards,
and revised or dropped reasons. Some evidence articles were initially
presented without later opposing-perspective paragraphs, which were restored
in the mixed-evidence phase. Coaching, external representation, evidence
access, competitive goals and peer response are substantive ingredients.
The final debate used rotating speakers and short team consultations.
None was independently varied to identify one unique active component
\cite[pp.~1250--1251]{KuhnUdell2003Development}.

Individual assessments before and after activity elicited the participant's
opinion and supporting reasons, then all reasons that someone might offer
on either side. Unique reasons indexed topic-specific knowledge. Product
quality used a hierarchy: unsupported or weakly justificatory reasons,
conditions/administrative objections, reasons addressing the institution's
functions, and functional reasons comparing it with alternatives. An
argument could also become two-sided by including reasons on both sides.
That is the authors' evaluative rubric; an administrative injustice does
not become morally unimportant simply because the rubric places it below
a comparative functional reason. Seventy percent of product responses were
double coded: 81 percent agreement, Cohen \(\kappa=.67\).

Discourse assessments paired students with opposing opinions, usually with
the same partner at pre/post assessment, for videotaped ten-minute dialogues.
Utterances were segmented and coded by their relation to the opponent's
preceding move. \emph{Counter-A} introduced an alternative argument;
\emph{Counter-C} criticized that move. \emph{Clarify} elaborated one's
own position; \emph{Add} elaborated conversationally without advancing
the argumentative goal. A rebuttal sequence was a counter followed directly
by Counter-C. The full source appendix distinguishes further questions,
agreement, interpretation and other functions. Off-topic and meta-dialogue
utterances were excluded from the proportion analyses. Training used
15 percent of dialogues, and a fourth rater recoded 25 percent of the
remaining coded material: 90 percent agreement, \(\kappa=.72\).
Coding reliability concerns the specified categories, not truth of the
participants' substantive claims
\cite[pp.~1248--1249,1252--1253,1259--1260]{KuhnUdell2003Development}.

For the statistics in these records, \(F\) is a reported analysis-of-variance
statistic, \(t\) a reported t-test statistic and \(p\) a statistical p-value.
One experimental dyad's recording was lost, leaving 32 participants for
discourse analysis, nineteen experimental and thirteen comparison. Mean
utterance counts per participant/dialogue changed from 7.23, SD~5.06,
to 29.58, SD~17.56, in the full condition, versus 7.15, SD~4.25,
to 11.11, SD~9.83, in comparison. These are within-dialogue utterances,
not independent participant recruitments. To separate composition from
volume, the authors analyzed utterance proportions with an arcsin
transformation and pretest covariates.

The overall group effect was \(F(5,25)=6.0,p=.001\). Counter-C
proportions rose from 5.3 to 30.6 percent in the full condition and
11.4 to 21.4 percent in comparison; its adjusted group test was
\(F(1,29)=4.21,p=.049\). Counter-A changed from 7.6 to 13.6
versus 5.3 to 1.2 percent, \(F(1,29)=21.07,p=.001\).
Clarify declined from 37.2 to 18.5 versus 44.3 to 32.4 percent.
Simple disagreement did not meet .05 in its group test, \(p=.08\).
The source's effects are therefore not uniform across discourse categories
\cite[p.~1252, tables~3--4]{KuhnUdell2003Development}.

At least one Counter-C occurred for all nineteen experimental discourse
participants at posttest, versus ten of thirteen comparison participants;
initial counts were eight and five. Rebuttal-sequence means rose from .16,
SD~.50, to 3.79, SD~1.75, in the full condition, versus .23,
SD~.44, to .46, SD~.66, in comparison. The time/condition interaction
was \(F(1,30)=32.37,p<.001\). The paper prints no confidence intervals
for these contrasts and reports no dyad/team/school-adjusted outcome model.
The authors also stress that advanced strategies remained uncommon:
clarification requests and more complex sequences did not show a broad
uniform advance. Observed Counter-C usage is not mastery of every
argumentative skill
\cite[pp.~1253,1256--1257, table~6]{KuhnUdell2003Development}.

Individual-product quality used all 34 completers: positive change was
reported for 15/21 full-condition participants and 5/13 comparison
participants. The article describes that between-condition difference as
not reaching statistical significance. Eight of the fifteen experimental
improvements included a shift to two-sided argument. The knowledge-of-reasons
analysis instead omitted one experimental participant with a missing
all-reasons response. Per-assessment reason counts rose from 3.30 to 4.45
in that group and 2.31 to 2.54 in comparison, but the between-group gain
difference was not significant. Distinct reasons pooled across the two
assessments averaged 6.70 versus 4.38, with \(t(31)=3.23,p=.003\).
That pooled repertoire statistic is not the posttest gain or a new sample
\cite[pp.~1254--1255, table~7]{KuhnUdell2003Development}.

The authors themselves retain the possibility that extending Phase~I
could improve the comparison group because its total exposure was shorter.
This limits a claim that interaction with an opposing side alone caused
the difference. The outcomes concern one practiced topic and its pre/post
dialogue/product tasks. No delayed novel-topic transfer test is reported
in this experiment; transfers cited in its discussion belong to other
studies. The experiment supports changes under the implemented package,
with attrition, small samples and dependence retained, not a universal
age transition or a demonstration that every debate improves rationality
\cite[pp.~1256--1258]{KuhnUdell2003Development}.

\subsection{Kuhn and Crowell's 2011 study: extended practice and novel essays}
\label{sec:p19-r151-kuhn-crowell}

The 2011 intervention tests a different transfer: from sustained dialogic
activity to individually written essays on assessment topics not taught in
that activity. Entering sixth graders were eleven or twelve and attended
one academically challenging urban public middle school in the northeastern
United States. The neighborhood was ethnically diverse and low-to-middle
income; the article reports 80 percent Hispanic or African-American and
60 percent free/reduced-lunch eligibility. Individual age ranges beyond
the entry description, native-language eligibility and a probability
recruitment frame are not supplied.

The school formed three classes of thirty to thirty-two students matched
on gender, ethnicity and standardized scores. Two classes were randomly
selected for intervention and one for comparison. This is random assignment
at the level of only three classes, not independent randomization of
each student's curriculum. Table~\ref{tab:p19-r151-kuhn-crowell-samples}
keeps the overlapping and distinct outcome datasets separate
\cite[pp.~546--547]{KuhnCrowell2011Dialogic}.

\begin{table}[htbp]
\centering\small
\begin{tabularx}{\textwidth}{@{}Xrr@{}}
\hline
Dataset & Intervention & Comparison\\
\hline
Main longitudinal sample & 48 & 23\\
Expanded main Year~3 posttest sample & 60 & 28\\
External Year~3 comparison school & --- & 50\\
Following-entry replication & 37 & 21\\
\hline
\end{tabularx}
\caption{Student counts for distinct outcome samples in Kuhn--Crowell 2011. The first two
rows overlap and must not be added as independent recruited populations;
the external group was matched rather than randomly assigned.}
\label{tab:p19-r151-kuhn-crowell-samples}
\end{table}

In the longitudinal sample there were 27 girls/21 boys in intervention and
11 girls/12 boys in comparison. Most supplied at least three of four
assessments; one comparison student supplied only initial/final assessments.
Three missed the final assessment and seven missed one interim assessment.
The article does not specify an outcome-level missing-data/imputation
procedure that reproduces all printed inferential denominators. The
expanded posttest samples included pupils absent initially or arriving
later, with sex counts 31/29 and 13/15. The external fifty eighth graders
came from another city school matched on scores, lunch eligibility and
ethnic composition, not from the intervention's randomization.

Both curricula met twice weekly for fifty-minute philosophy classes.
The intervention spent about thirteen sessions per topic, four topics
per year, across three years. Same-side preparation generated and ranked
reasons. Partners then conducted electronic dialogues with six opposing
pairs, preserving chat transcripts for reflection. Group preparations,
a rotating-speaker debate, argument-map debriefing and a final essay
completed a topic cycle. Evidence materials were added from the final
topic of Year~1; students increasingly generated questions themselves,
with coaches supplying concise answers. Reflection changed form in
Year~3, and topics became three-sided. Note~2 records differences
between the two classes' Year~1 time allocation. This is a multicomponent,
changing curriculum rather than one constant dialogue-dose treatment.

Comparison students received teacher-led whole-class discussion, some
activities such as dramatization and more frequent essays: fourteen per
year versus four in intervention. Thus the control was educationally
active and had more writing practice, while coaching, peer structure,
reflection and evidence access differed. It does not isolate the causal
contribution of electronic rather than spoken dialogue or any single
ingredient
\cite[pp.~547--548,552, note~2]{KuhnCrowell2011Dialogic}.

Teacher-pay essays were assessed initially and after each of three years.
A separate assisted-dying essay appeared only at the Year~3 endpoint.
These were assessment subjects, not the curriculum's deeply studied
topics; neither political nor medical position was the answer key.
At Year~3, essays received twenty minutes and follow-up questions ten
minutes, with the teacher-pay assessment first on a different day.
The follow-up asked which unanswered questions would help the argument;
those assessment questions were not answered. More questions therefore
measure proposed evidence seeking, not actual verified information use.

The scoring categories preserve the source's distinctions. An own-side
argument offers advantages of the favored position. A \emph{dual perspective}
also attacks the opposing option, so all considerations can still favor
the same conclusion. An \emph{integrative perspective} acknowledges a
disadvantage of the favored position or advantage of its opponent, requiring
considerations leading in opposite directions to be addressed. This is
not identical to the 2003 paper's one-/two-sided argument coding.
Idea-unit coding was blind to condition and time; a random third of
essays was double coded. Agreement was 88 percent, \(\kappa=.76\),
for teacher pay and 93 percent, \(\kappa=.91\), for assisted dying.
Reliability of classification does not itself validate every factual
premise or normative recommendation within an essay
\cite[pp.~546,548--550, table~1]{KuhnCrowell2011Dialogic}.

For teacher-pay dual-perspective counts, the reported condition/time
interaction was \(F(3,67)=6.11,p<.001,\eta_p^2=.216\).
At Year~3, the percentages producing at least one dual-perspective
argument were 79 versus 29; integrative arguments appeared in 30 versus
zero percent, having been absent at earlier main-study assessments.
Both groups' essays became longer. Year~3 total argument means were
3.96 versus 2.61; the source also reports dual-perspective shares .49
versus .24 after accounting for overall production. Figure~1 error bars
are standard errors, not 95 percent confidence intervals.

On the new assisted-dying topic, dual-perspective argument counts averaged
1.17 versus .36, \(t(86)=3.85,p<.001\). Percentages supplying at
least one dual-perspective argument were 73, 29 and 24 for intervention,
internal comparison and external comparison; integrative percentages were
30, 18 and eight. These are different endpoints with different samples,
not universal success percentages. Proposed question counts averaged
4.18 versus .33 for teacher pay and 3.26 versus 1.64 for assisted dying,
with .92 in the external comparison. Counts and general-versus-case-specific
question coding support a limited evidence-seeking interpretation; the
article does not establish the future truth or usefulness of every question
\cite[pp.~549--551, tables~2--5]{KuhnCrowell2011Dialogic}.

The next incoming cohort provided a within-school replication through
Year~2, not an independent multi-population replication. After classroom
reassignment for behavior-management reasons, all students switching
conditions were excluded; note~4 does not give that exclusion count.
The retained sample was 37 intervention and 21 comparison students,
with comparison assessments available only at Year~2. Two intervention
responses were missing initially and two at Year~1. At Year~2,
dual-perspective count means were 1.05 versus .43,
\(t(56)=2.81,p=.007\), and integrative arguments appeared in 14
versus zero percent. Proposed question means were 1.48 versus .90,
but their difference was not significant, \(p=.148\). Replication
therefore does not mean every main-study outcome repeated at the same
strength or timing
\cite[pp.~550--552, table~6, note~4]{KuhnCrowell2011Dialogic}.

The paper supplies no interval estimates for these principal contrasts
and no reported class-cluster-adjusted model. Treating every pupil as
an independently randomized intervention unit would overstate the
randomization evidence. Attrition, changing membership, excluded
class-switchers, differences in adult facilitation and a restricted
school setting qualify causal and generalization claims. Nevertheless,
novel-topic individual essays provide an actual transfer outcome absent
from the undergraduate selection study and the same-topic 2003 experiment.
They do not establish transfer to abstract deduction, measured working-memory
capacity, an explicit concept of logical validity, all populations or
long-term performance after the curriculum ends.

\subsection{What these comparisons leave distinct}
\label{sec:p19-r151-social-capacity-boundary}

A complete selection under a stipulated conditional rule, a partner-responsive
counterargument, a higher-rated individual argument and an integrative essay
are different achievements. The first permits an exact truth-preservation
benchmark; the others use declared task-specific functional or evaluative
rubrics. The studies supply evidence at their measured times and scopes,
not one shared rationality score. The capacity and metalogical accounts in
this chapter additionally concern resource limits and explicit understanding.
These social-intervention outcomes do not directly measure either construct
or isolate capacity growth as their causal mechanism. A developmental
comparison must preserve what was practiced, what was tested, who was
sampled and whether a new topic or later occasion was actually observed.

